\documentclass{article}
\usepackage{amsmath,amssymb}
\usepackage{xurl}
\usepackage[hidelinks]{hyperref}
\setlength{\emergencystretch}{2em}
% Shared commands for notes in docs/paper-gaps/.

\DeclareUnicodeCharacter{2080}{\ifmmode{}_{0}\else\textsubscript{0}\fi}
\DeclareUnicodeCharacter{2081}{\ifmmode{}_{1}\else\textsubscript{1}\fi}
\DeclareUnicodeCharacter{2082}{\ifmmode{}_{2}\else\textsubscript{2}\fi}
\DeclareUnicodeCharacter{2099}{\ifmmode{}_{n}\else\textsubscript{n}\fi}

\newcommand{\Lean}{\textsc{Lean}}
\ExplSyntaxOn
\NewDocumentCommand{\leanid}{m}{
  \ifmmode
    \textsf{\detokenize{#1}}
  \else
    \group_begin:
    \urlstyle{sf}
    \tl_set:Nn \l_tmpa_tl {#1}
    \tl_replace_all:Nnn \l_tmpa_tl {\_} {_}
    \exp_args:NV \path \l_tmpa_tl
    \group_end:
  \fi
}
\ExplSyntaxOff
\providecommand{\tr}{\operatorname{tr}}
\newcommand{\tnleanIssueBase}{https://github.com/LionSR/TNLean/issues/}
\newcommand{\tnleanPrBase}{https://github.com/LionSR/TNLean/pull/}
\newcommand{\tnleanDocBase}{https://sirui-lu.com/TNLean/docs/}
\newcommand{\ghissue}[1]{\href{\tnleanIssueBase#1}{GitHub issue \##1}}
\newcommand{\ghpr}[1]{\href{\tnleanPrBase#1}{PR \##1}}
\newcommand{\doclink}[1]{\href{\tnleanDocBase#1.html}{\path{#1}}}

% Dirac notation and the identity operator, as in blueprint/src/macros/common.tex.
% \providecommand keeps note-local definitions authoritative.
\usepackage{dsfont}
\providecommand{\ket}[1]{|#1\rangle}
\providecommand{\bra}[1]{\langle #1|}
\providecommand{\braket}[2]{\langle #1 | #2 \rangle}
\providecommand{\ketbra}[2]{|#1\rangle\!\langle #2|}
\providecommand{\Id}{\mathds{1}}
\providecommand{\one}{\mathds{1}}
\providecommand{\C}{\mathbb{C}}
\providecommand{\MN}[1]{M_{#1}(\C)}
\providecommand{\MD}{M_D(\C)}

% External referencing for paper sources.  The build helper writes sanitized
% auxiliary files under build/paper-aux/ and excludes duplicated source labels.
\usepackage{xr}

% Import a generated paper auxiliary file with a caller-supplied label prefix.
% The candidate paths support builds from docs/paper-gaps/, the repository root,
% and build/paper-aux/ itself.
\newcommand{\importpaperaux}[2]{%
  \IfFileExists{../../build/paper-aux/#2.aux}{%
    \externaldocument[#1]{../../build/paper-aux/#2}%
  }{%
    \IfFileExists{build/paper-aux/#2.aux}{%
      \externaldocument[#1]{build/paper-aux/#2}%
    }{%
      \IfFileExists{#2.aux}{%
        \externaldocument[#1]{#2}%
      }{}%
    }%
  }%
}

% General external reference macros
\newcommand{\paperref}[2]{\ref{#1-#2}}
\newcommand{\papereqref}[2]{(\ref{#1-#2})}

% CPSV16 source (1606.00608 / MPDO-22-12-17-2.tex) external document and shortcuts
\importpaperaux{cpsv16-}{1606.00608/MPDO-22-12-17-2-tnlean}
\newcommand{\cpsvref}[1]{\paperref{cpsv16}{#1}}
\newcommand{\cpsveqref}[1]{\papereqref{cpsv16}{#1}}

% Verdict marker: \gapnote{<kind>}{<status>}. Kind is the policy
% classification (clarification, local-correction, scope-restriction,
% unfaithful, false-source, open-gap); status is open, wip (elimination
% underway), resolved, or historical. Machine-read by the site index;
% prints a verdict line.
\newcommand{\gapnote}[2]{\par\noindent\textbf{Verdict:} #1 (#2).\par}

\title{Simultaneous Left Inverse of Injective Operator Blocks}
\author{}
\date{2026-09-25}
\begin{document}
\maketitle
\gapnote{local-correction}{open}

\section{Source statement}
To extract the $F$-symbols, Bultinck et al.\ apply a left inverse of the
labelled block tensors to both sides of an identity between sums over the
final label $d$~\cite{Bultinck2017Anyons}. They write that injectivity of the
single-block tensors implies that the matrix $(B_d^{ij})_{\alpha\beta}$, with
rows indexed by the physical pair $ij$ and columns by the bond pair
$\alpha\beta$, has a left inverse $B_d^+$ with
\begin{align}
    B_d^+B_{d'} & =\delta_{dd'}\,\mathbb 1_{\chi_d}\otimes\mathbb 1_{\chi_d}.
    \label{eq:bmwshv17_joint_inverse}
\end{align}
\footnote{Source traceability: \path{References/1511.08090/AnyonsPEPS.tex},
    line 269, within the derivation at lines 252--277.}
The case $d=d'$ is injectivity of $B_d$. The case $d\ne d'$ asks that one
left inverse for $B_d$ annihilates every other block.

\section{Counterexample}
Relation (\ref{eq:bmwshv17_joint_inverse}) for all labels together says
that the one-site map
\begin{align}
    \Phi:\mathbb C^{p^2}\longrightarrow\bigoplus_d M_{\chi_d},
    \qquad
    v\longmapsto\Bigl(\sum_{i,j}v_{ij}B_d^{ij}\Bigr)_d,
    \label{eq:bmwshv17_joint_map}
\end{align}
is surjective. The row of $(B_d^+)$ indexed by $(\alpha,\beta)$ is a preimage
of the matrix unit $E_{\alpha\beta}$ in the summand $d$. Surjectivity forces
\begin{align}
    p^2 & \ge\sum_d\chi_d^2 .
    \label{eq:bmwshv17_dimension_count}
\end{align}
Injectivity of each block separately does not imply it. Take physical
dimension $p=2$. Let $B_1$ have bond dimension $\chi_1=2$ and let its four
letters $B_1^{ij}$ form a basis of $M_2$. Let $B_2$ have bond dimension
$\chi_2=1$ with some nonzero letter. Both blocks are injective. They are not
gauge equivalent, since their bond dimensions differ. The map $\Phi$ goes from
$\mathbb C^4$ to $M_2\oplus M_1\cong\mathbb C^5$, so its rank is at most
$4<5$. Hence no simultaneous left inverse exists. The data are
nondegenerate: every block is injective and every weight is nonzero.

\section{Missing hypothesis}
The step needs joint linear independence of the labelled blocks at one site,
that is, surjectivity of (\ref{eq:bmwshv17_joint_map}). Injectivity of each
block together with pairwise non-equivalence does not supply it.

The formal complete zipper fusion family records the simultaneous left inverse
as a field, which is exactly this hypothesis. The construction of such a family
from split compressions of the pairwise products takes it as an input.
\footnote{Formal definitions: \leanid{MPOTensor.CompleteZipperFusionFamily}
    and \leanid{MPOTensor.CompleteZipperFusionFamily.ofCompression}.}

\section{Elimination route}
Blocking should remove the hypothesis. For pairwise non-equivalent normal
tensors, the words of a common sufficiently large length $L$ are expected to
span the full direct sum $\bigoplus_d M_{\chi_d}$, by the linear independence
of non-equivalent normal tensors. The blocked tensors would then satisfy
(\ref{eq:bmwshv17_joint_inverse}) at one blocked site, and the $F$-symbol
extraction would run at length $L$. This route is plausible but is not proved
here. The zipper and reconstruction identities are also needed at the blocked
level, where they follow from the one-site identities by composing fusion
tensors along the word.

\bibliographystyle{alpha}
\bibliography{references}
\end{document}
